% SEGMENT OLP-0555-S01
\documentclass[../../../include/open-logic-section]{subfiles}

\begin{document}
\olfileid{sth}{ordinals}{zfm}	
\olsection{$\ZFminus$: ஒரு முக்கியக் கட்டம்}

மாற்றீட்டை எவ்வாறு நியாயப்படுத்துவது (அப்படி நியாயப்படுத்த முடிந்தால்) என்ற
கேள்விக்கு எளிய விடை இல்லை. எனவே அதை \olref[replacement][]{chap} க்காக
ஒதுக்கிவைப்போம். ஆனால் மாற்றீட்டைச் சேர்த்ததால், இன்னொரு முக்கியக் கட்டத்தை
அடைந்துள்ளோம். $\ZFminus$ என்ற கோட்பாட்டுக்குத் தேவையான எல்லா அடிகோள்களும்
இப்போது நம்மிடம் உள்ளன. விரிவாக:

\begin{defn}
$\ZFminus$ கோட்பாடு பின்வரும் அடிகோள்களைக் கொண்டுள்ளது: விரிவுநிலை,
சேர்ப்பு, ஜோடிகள், அடுக்குக்கணங்கள், முடிவிலி ஆகிய அடிகோள்களும், பிரிப்பு
மற்றும் மாற்றீட்டுத் திட்டங்களின் எல்லா நேர்வுகளும். வேறுவிதமாகச் சொன்னால்,
$\ZFminus$ ஆனது $\Zminus$ இல் மாற்றீட்டைச் சேர்க்கிறது.
\end{defn}
\noindent
இது \emph{செர்மெலோ--ஃபிராங்கல்} கணக் கோட்பாட்டைக் குறிக்கிறது (பின்னர்
காணவிருக்கும் ஏதோ ஒன்றை \emph{நீக்கிய} வடிவம்). \citeyear{Fraenkel1922}
இல் மாற்றீட்டை வடிவமைத்த பெருமை ஃபிராங்கலுக்குக் கொடுக்கப்படுவதால் அவரது
பெயர் இடம்பெறுகிறது; எனினும் முதல் துல்லியமான வடிவமைப்பு
\citet{Skolem1922} ஆல் வழங்கப்பட்டது.

\end{document}
